% ===================================================================
%  పట్టిక సంఖ్య 3 — బాల్యం మరియు యౌవనం.
%
%  Traduction, faite en 2026, en telougou d'aujourd'hui, sur l'ido de
%  texto/io. Regles de la colonne : voir 10-pattika-01.tex. Deux
%  scenes, deux numerotations : les cles portent c1 et c2.
%
%  LA PARENTE OBLIGE LE TELOUGOU A CHOISIR LA OU L'IDO RESTE VAGUE,
%  ET C'EST LE MEME EMBARRAS QU'AU TABLEAU 4 MARATHE. L'ido dit
%  onklulo, onklino, kuzulo, kuzino sans dire de quel cote ; le
%  telougou a un mot par relation et il faut lire la planche pour
%  savoir lequel. On ecrit మామయ్య pour l'oncle et అత్తయ్య pour la
%  tante — les termes que l'enfant emploie dans la vie —, చిన్నాన్న
%  కొడుకు et చిన్నాన్న కూతురు pour les cousins du cote paternel, et
%  అన్న / చెల్లెలు, qui distinguent l'aine du cadet la ou le francais
%  ne distingue que le sexe.
%
%  LE PRINTEMPS ET LA FETE SONT ENTIEREMENT TELOUGOUS :
%
%    వసంతకాలం   le printemps    వేసవి      l'ete
%    మబ్బు      le nuage        ఆకాశం      le ciel
%    మొగ్గలు    les bourgeons   పూలు       les fleurs
%    తేనెటీగ    l'abeille       తేనెతుట్టె  la ruche
%    తేనె       le miel         గాడిద      l'ane
%    సీతాకోకచిలుక le papillon    గులాబీ     la rose
%    ముల్లు     l'aiguille de l'horloge
%    శిలువ      la croix        గంట        la cloche
%    ఊతకర్ర     la bequille     బిచ్చగాడు  le mendiant
%    రెల్లు     les roseaux     బాణసంచా    le feu d'artifice
%    పడవ        la barque       తెడ్డు     la rame
%    చుక్కాని   le gouvernail   రేవు       le quai
%    వంతెన      le pont         ఒడ్డు      la rive
%    తిరునాళ్ళు  la fete foraine ఎద్దు      le taureau
%    పులి       le tigre        సింహం      le lion
%    జూలు       la criniere     ఏనుగు      l'elephant
%    తొండం      la trompe       మరుగుజ్జు  le nain
%
%  ET « ఊయల » EST A LA FOIS LE BERCEAU ET LA BALANCOIRE, comme au
%  tableau 2. Ce n'est pas une commodite : les deux se suspendent et
%  se balancent, et le telougou les a nommes d'un seul mot — la
%  quatrieme homonymie vieille du releve, apres నాంగర్ / नांगर,
%  घाट et बंब en marathe.
%
%  UN EMPRUNT QUE JE NE SAIS PAS EVITER : l'hirondelle. Le telougou
%  a des noms pour la plupart des oiseaux de ce livret, mais je n'ai
%  pas de forme sure pour celui-la, et en fabriquer une serait pire
%  que d'emprunter. On ecrit స్వాలో పక్షులు, et le trou est note ici
%  plutot que masque.
%
%  DEUX INVERSIONS PREVENUES — la fleche du clocher et les abeilles
%  de la ruche —, ET NEUF QUI ONT PASSE. Le compte est mauvais, et il
%  merite d'etre compris plutot que deplore.
%
%  LE TELOUGOU MORD PLUS FORT QUE LE MARATHI, ET POUR UNE AUTRE
%  RAISON. Le marathi postposait : « X de Y » sortait a l'envers, et
%  le remede etait de couper la phrase en deux. Le telougou, lui,
%  AGGLUTINE : « la laitiere avec son ane » se dit naturellement
%  « తన గాడిదతో పాలమ్మే ఆమె », et l'ane entre DANS le groupe nominal
%  de la laitiere avant qu'elle soit nommee. Ce n'est plus un ordre
%  de mots qu'on redresse, c'est un groupe qu'il faut ouvrir. Le
%  remede est donc un cran plus lourd : detacher le second terme en
%  apposition, apres le premier, avec sa propre marque de cas.
%
%  C'est pourquoi neuf blocs de ce tableau se lisaient parfaitement
%  et sortaient faux : le nuage avant le ciel, l'ane avant la
%  laitiere, la bequille avant le mendiant, le baton avant la vieille
%  femme, les musiciens avant les gamins. Aucune n'etait visible a
%  l'oeil ; renvoji.py les a toutes vues.
% ===================================================================

%%K t03-apar-1 apar
\VUsaut{3.0mm}
\VUtitre{15.0pt}{60}{పట్టిక సంఖ్య 3}
\VUsaut{1.6mm}
\VUtitre{11.4pt}{0}{బాల్యం మరియు యౌవనం.}
\VUsaut{3.4mm}

%%K t03-apar-2 apar
(3వ మరియు 4వ పట్టికలు, నాలుగు \VUgras{కుటుంబ రంగాల్లో} \VUgras{వాతావరణం}, \VUgras{వయసు},
\VUgras{కుటుంబం}, \VUgras{దుస్తులు} మరియు \VUgras{స్థితి} గురించిన ముఖ్యమైన భావనలను చూపేలా
కలిపి కూర్చబడ్డాయి.)

%%K t03-tit-1 sub
\VUcentre{12.0pt}{0}{\textit{మొదటి రంగం.}}
\VUsaut{3.0mm}

%%K t03-tit-2 sub
\VUpk{10.2pt}{40}{పల్లెలో బాప్తిస్మ వేడుక.}
\VUsaut{2.4mm}

%%K t03-tit-3 sub
\VUpk{10.2pt}{40}{వసంతకాలం.}
\VUsaut{3.0mm}

%%K t03-c1-01-1 p
1. --- మనం \VUgras{వసంతకాలంలో} ఉన్నాం — \VUgras{మార్చి}, \VUgras{ఏప్రిల్} లేదా \VUgras{మే}
నెలల్లో, \VUgras{వారంలోని} \VUgras{సోమవారం}, \VUgras{మంగళవారం} లేదా \VUgras{బుధవారం} రోజుల్లో.
ఉదయం \VUgras{పది గంటలు}.

%%K t03-c1-02-1 p
2. --- \VUgras{వాతావరణం} బాగుంది, \VUgras{గాలి} స్వచ్ఛంగా ఉంది. సూర్యుడు ప్రకాశిస్తున్నాడు.
కొన్ని చిన్న \VUgras{మబ్బులు} \textsuperscript{(2)} పైకి లేస్తున్నాయి, నీలి \VUgras{ఆకాశంలో}
\textsuperscript{(1)}.

%%K t03-c1-03-1 p
3. --- \VUgras{చెట్లకు} ఇంకా \VUgras{ఆకులు} అంతగా రాలేదు. సన్నని \VUgras{చినార చెట్లకు}
\textsuperscript{(3)} మరియు ఎత్తైన \VUgras{ప్లేన్ చెట్టుకు} \textsuperscript{(17)} ఇప్పటివరకు
\VUgras{మొగ్గలు} మాత్రమే ఉన్నాయి.

%%K t03-c1-04-1 p
4. --- \VUgras{స్వాలో పక్షులు} \textsuperscript{(4)} తిరిగి వచ్చి, ఆనందంగా కిచకిచలాడుతూ ఆ
\VUgras{బాణం} \textsuperscript{(7)} చుట్టూ ఎగురుతున్నాయి — \VUgras{గంటస్తంభం}
\textsuperscript{(5)} పైనున్న బాణం, ఆ స్తంభం పల్లె \VUgras{చర్చికి} \textsuperscript{(6)}
కిరీటంలా ఉంది.

%%K t03-tit-4 sub
\VUsaut{3.4mm}
\VUpk{10.2pt}{40}{చర్చి.}
\VUsaut{3.0mm}

%%K t03-c1-05-1 p
5. --- తన పెద్ద \VUgras{ప్రధాన ద్వారంతో} \textsuperscript{(13)}, లావు \VUgras{గడియారంతో}
\textsuperscript{(8)} ఆ చర్చి చాలా సాదాసీదాగా ఉంది. చిన్న \VUgras{ముల్లు}
\textsuperscript{(9)} X అంకెపై ఉంది, అది \VUgras{గంటలను} చూపుతుంది. \VUgras{పెద్దది}
\textsuperscript{(10)} XII పై ఉంది (మధ్యాహ్నం); అది \VUgras{నిమిషాలను} చూపుతుంది.

%%K t03-c1-06-1 p
6. --- గంటస్తంభంలో \VUgras{గంటలు} \textsuperscript{(11)} ఆనందంగా మోగుతున్నాయి. పైకప్పు శిఖరాన
ఒక చిన్న రాతి \VUgras{శిలువ} \textsuperscript{(12)} నిలబడి ఉంది.

%%K t03-c1-07-1 p
7. --- చర్చికి పెద్ద \VUgras{ప్రధాన ద్వారం} \textsuperscript{(13)} మాత్రమే కాదు, పక్కన ఒక
చిన్న తలుపు కూడా ఉంది. ఆ చిన్న తలుపు గుండా \VUgras{మతగురువు} \textsuperscript{(15)} తన
\VUgras{సుతానా} వేసుకుని \VUgras{పూజాసామగ్రి గదిలోంచి} \textsuperscript{(14)} బయటికి వచ్చి
\VUgras{తన ఇంటికి} \textsuperscript{(19)} వెళ్తున్నారు.

%%K t03-c1-08-1 p
8. --- ఆయన దగ్గరే ఇదిగో \VUgras{పాలమ్మే ఆమె} \textsuperscript{(24)}, తన \VUgras{గాడిదతో}
\textsuperscript{(23)}. పిల్లలకు మంచి పాలు తీసుకెళ్ళిన ఊరి నుంచి ఆమె తిరిగి వస్తోంది.

%%K t03-tit-5 sub
\VUsaut{4.0mm}
\VUpk{10.2pt}{40}{పండ్ల తోట.}
\VUsaut{3.4mm}

%%K t03-c1-09-1 p
9. --- ఆ ఉదయపు నడక గురించి అలిబోరాన్ గారికి (గాడిదకు) పూర్తి తృప్తి. \VUgras{పూల}
\textsuperscript{(20)} సువాసనతో నిండిన గాలిని అది ఆనందంగా పీలుస్తోంది — ఆ పూలు పక్కనున్న
\VUgras{పండ్ల తోటలోని} \VUgras{పండ్ల చెట్లను} \textsuperscript{(18)} కప్పుతున్నాయి. ఆ
\VUgras{పీచు చెట్లు} \textsuperscript{(18)} మరియు \VUgras{ఆప్రికాట్ చెట్లు}
\textsuperscript{(19)} పూర్తిగా గులాబీ, తెలుపు రంగుల్లో ఉన్నాయి, మంచి పంట వస్తుందనే ఆశ
కలిగిస్తున్నాయి.

%%K t03-c1-10-1 p
10. --- ఈలోగా చిన్న \VUgras{తేనెటీగలు} \textsuperscript{(22)} — \VUgras{తేనెతుట్టెలోని}
\textsuperscript{(21)} తేనెటీగలు — ఝంకారం చేస్తూ తమ తియ్యని \VUgras{తేనెను} పోగుచేసుకునేందుకు
అక్కడికి వెళ్తున్నాయి.

%%K t03-c1-11-1 p
11. --- కానీ ఏమవుతోంది? ఇదిగో పల్లె పిల్లలు పరిగెత్తుకుంటూ వచ్చి భయపడిన \VUgras{బాతులను}
\textsuperscript{(25)} పారదోలుతున్నారు. అది నోటరీ కూతురి \VUgras{బాప్తిస్మం} తర్వాత చర్చి
నుంచి ఊరేగింపు బయటికి రావడం.

%%K t03-c1-12-1 p
12. --- తన \VUgras{ఊయలలో} \textsuperscript{(26)} ఉన్న చిన్న జేకబ్ గంటల ఆనందపు మోతకు
మేల్కొంటున్నాడు, ఆ బాప్తిస్మ ఊరేగింపును తానూ చూడాలనుకున్న అతని అక్క మాగ్దలేనా, ఇంటి గుమ్మం
దగ్గర తన జుట్టు దువ్వి బట్టలు వేసేందుకు రమ్మని తల్లిని బతిమాలుతోంది.

%%K t03-c1-13-1 p
13. --- తల్లి ఒప్పుకుంటుంది. ఆమె తన \VUgras{తలవస్త్రం} \textsuperscript{(28)}, తన
\VUgras{రవిక} \textsuperscript{(29)} మరియు తన \VUgras{నడుము గుడ్డ} \textsuperscript{(30)}
వేసుకుని తలుపు ముందున్న చిన్న కుర్చీపై కూర్చుంటుంది. మాగ్దలేనాకు తన \VUgras{లోపలి లంగా}
\textsuperscript{(31)} మరియు తన \VUgras{కార్సెట్} \textsuperscript{(32)} మాత్రమే ఉన్నాయి.

%%K t03-c1-14-1 p
14. --- ఆమె ఎంత సంతోషంగా ఉంది! తనకిష్టమైన పూలను ఆమె మరిచిపోతోంది — తన \VUgras{కార్నేషన్లు}
\textsuperscript{(34)}, వాటిపై వాలే \VUgras{సీతాకోకచిలుకలు} \textsuperscript{(35)}, తన
\VUgras{తులిప్‌లు} \textsuperscript{(36)}, తన \VUgras{లోయ లిల్లీలు} \textsuperscript{(37)},
\VUgras{కుండీల్లో} \textsuperscript{(38)}, \VUgras{గోడపై} \textsuperscript{(33)}, ఇంకా అంత
అందమైన కంచెను ఏర్పరిచే తన \VUgras{గులాబీలు} \textsuperscript{(39)}. ఊరేగింపులోని స్త్రీల
ఖరీదైన దుస్తులనే ఆమె చూస్తోంది.

%%K t03-c1-15-1 p
15. --- పల్లె పిల్లలు దుస్తులను పెద్దగా పట్టించుకోరు. \VUgras{మిఠాయిలు} \textsuperscript{(55)}
లేదా \VUgras{నాణేలు} \textsuperscript{(52)} దక్కించుకునేందుకు వాళ్ళు ఒకరినొకరు తోసుకుంటూ
ముందుకు దూసుకుపోతున్నారు. వాళ్ళలో ఒకడు \VUgras{ఏడుస్తున్నాడు} \textsuperscript{(40)}.

%%K t03-c1-16-1 p
16. --- ఇంకొకడు \VUgras{కుక్క} \textsuperscript{(42)} కాలు తొక్కాడు, అది అరుస్తూ పారిపోయి
\VUgras{కోడిని} \textsuperscript{(43)} మరియు దాని \VUgras{కోడిపిల్లలను} \textsuperscript{(44)}
పారదోలుతోంది.

%%K t03-tit-6 sub
\VUsaut{5.0mm}
\VUpk{10.2pt}{40}{బాప్తిస్మ ఊరేగింపు.}
\VUsaut{4.0mm}

%%K t03-c1-17-1 p
17. --- ఊరేగింపుకు ముందు \VUgras{దాది} \textsuperscript{(45)} నడుస్తోంది. ఆమెకు పొడవాటి
రిబ్బన్లతో అందమైన \VUgras{తలకట్టు} \textsuperscript{(46)} ఉంది. \VUgras{నవజాత శిశువును}
\textsuperscript{(47)} ఆమె ఎత్తుకుని ఉంది, ఆ బిడ్డ \VUgras{లేసులతో} \textsuperscript{(48)}
పూర్తిగా అలంకరించి ఉంది.

%%K t03-c1-18-1 p
18. --- దాది వెనుక \VUgras{బాప్తిస్మ తండ్రి} \textsuperscript{(49)} మరియు
\VUgras{బాప్తిస్మ తల్లి} \textsuperscript{(53)} వస్తున్నారు. మిఠాయిలూ నాణేలూ వాళ్ళే
పంచుతున్నారు. బాప్తిస్మ తండ్రికి \VUgras{చేతి తొడుగులు} \textsuperscript{(51)} మరియు
\VUgras{సిలిండర్ టోపీ} \textsuperscript{(50)} ఉన్నాయి. బాప్తిస్మ తల్లి చేతిలో అందమైన
\VUgras{పూలగుత్తి} \textsuperscript{(54)} ఉంది.

%%K t03-c1-19-1 p
19. --- వాళ్ళ వెనుక పిల్లవాడి \VUgras{మామయ్య} \textsuperscript{(56)} మరియు \VUgras{అత్తయ్య}
\textsuperscript{(58)} మనకు కనిపిస్తారు. అత్తయ్యకు హుందా అయిన \VUgras{టోపీ}
\textsuperscript{(59)}, మామయ్యకు నల్ల \VUgras{రెడింగోట్} \textsuperscript{(57)} మరియు తెల్ల
చొక్కా ఉన్నాయి. ఆ తర్వాత ఇదిగో \VUgras{చిన్నాన్న కొడుకు} \textsuperscript{(60)} మరియు
\VUgras{చిన్నాన్న కూతురు} \textsuperscript{(61)}. ఆమె నవజాత శిశువు \VUgras{చెల్లెలైన}
\textsuperscript{(62)} చిన్న బెర్తా చేయి పట్టుకుని ఉంది.

%%K t03-c1-20-1 p
20. --- బెర్తా ఇంకో \VUgras{అన్న} \textsuperscript{(63)} వెనుక నడుస్తున్నాడు. అతను ఒక
\VUgras{బంధువుకు} \textsuperscript{(64)} చేయి ఇస్తున్నాడు. కుటుంబ \VUgras{స్నేహితులు}
\textsuperscript{(65)} ఆ తర్వాత వస్తున్నారు.

%%K t03-tit-7 sub
\VUsaut{4.6mm}
\VUpk{10.2pt}{40}{చూసేవాళ్ళు.}
\VUsaut{3.4mm}

%%K t03-c1-21-1 p
21. --- ఊరేగింపు దగ్గర ఇదిగో ఒక \VUgras{బిచ్చగాడు} \textsuperscript{(66)}, తన
\VUgras{ఊతకర్రపై} \textsuperscript{(67)} ఆనుకుని. అతను \VUgras{భిక్ష} అడుగుతున్నాడు.

%%K t03-c1-22-1 p
22. --- మరోవైపు చాలా \VUgras{మర్యాదస్థుడైన} \textsuperscript{(68)} ఒక చిన్న అబ్బాయి ఉన్నాడు.
తనకు తెలిసినవాళ్ళకు అతను నమస్కరించి « \VUgras{నమస్కారం} » అంటున్నాడు.

%%K t03-c1-23-1 p
23. --- బాప్తిస్మ ఊరేగింపును చూసేందుకు పల్లెవాసులు తమ ఇళ్ళలోంచి బయటికి వచ్చారు. తన
\VUgras{నూలు టోపీతో} ఒక \VUgras{పల్లెవాసి} \textsuperscript{(69)}, పక్కనే ఒక
\VUgras{పల్లెవాసిని} \textsuperscript{(70)} మనకు కనిపిస్తారు. ఆ తర్వాత ఒక \VUgras{ముసలామె}
\textsuperscript{(71)}, తన \VUgras{కర్రపై} \textsuperscript{(72)} ఆనుకుని. చివరగా
\VUgras{మిల్లువాడు} \textsuperscript{(73)}, తన పెద్ద \VUgras{ఉన్ని టోపీతో}
\textsuperscript{(74)}, మరియు \VUgras{చెక్క చెప్పులు} \textsuperscript{(75)} వేసుకుని తన
పసిబిడ్డకు నడక నేర్పుతున్న ఒక యువ పల్లెవాసిని.

%%K t03-c1-24-1 p
24. --- ఆ రంగం చాలా సరదాగా ఉంది.

%%K t03-tit-8 sub
\VUcentre{12.0pt}{0}{\textit{రెండవ రంగం.}}
\VUsaut{3.0mm}

%%K t03-tit-9 sub
\VUpk{10.2pt}{40}{బహిరంగ ఉత్సవం.}
\VUsaut{2.4mm}

%%K t03-tit-10 sub
\VUpk{10.2pt}{40}{నౌకా క్రీడలు మరియు పడవ పందేలు.}
\VUsaut{3.0mm}

%%K t03-c2-01-1 p
1. --- \VUgras{వేసవి} వచ్చింది. \VUgras{జూన్} (జూలై లేదా \VUgras{ఆగస్టు}) నెలలోని మండే ఎండ
యువకులను స్నానానికీ పడవ నడపడానికీ పురిగొల్పుతోంది.

%%K t03-c2-02-1 p
2. --- నదికి కుడి ఒడ్డున, నౌకా క్రీడల్లోనూ పడవ పందేల్లోనూ పాల్గొనేందుకు పడవవాళ్ళు బట్టలు
విప్పుతున్నారు.

%%K t03-c2-03-1 p
3. --- వాళ్ళలో ఒకడు తన \VUgras{చెప్పులు} \textsuperscript{(1)} విప్పుతున్నాడు; వాటి కట్లు
(గుండీలు) విప్పుతున్నాడు. అతని ఇద్దరు స్నేహితులు ఇప్పటికే సిద్ధంగా ఉన్నారు. ఇప్పుడు వాళ్ళకు తమ
\VUgras{నేసిన చొక్కా} \textsuperscript{(2)} మరియు \VUgras{ఈత లాగు} \textsuperscript{(3)}
మాత్రమే ఉన్నాయి. స్నేహితుల కోసం ఎదురుచూస్తూ వాళ్ళు మాట్లాడుకుంటున్నారు. ఇంకో ఒడ్డున జరిగే
తిరునాళ్ళ గురించీ ఉత్సవం గురించీ వాళ్ళు మాట్లాడుతున్నారు.

%%K t03-c2-04-1 p
4. --- వాళ్ళ దగ్గర నేలపై స్నేహితుల \VUgras{బట్టలు} పడి ఉన్నాయి : ఒక జత \VUgras{బూట్లు}
\textsuperscript{(4)}, \VUgras{చెప్పులు} \textsuperscript{(5)}, \VUgras{మేజోళ్ళు}
\textsuperscript{(6)}, \VUgras{ఉన్ని} \textsuperscript{(7)} మరియు \VUgras{గడ్డి}
\textsuperscript{(8)} టోపీలు, ఒక \VUgras{మెడకట్టు} \textsuperscript{(9)}.

%%K t03-c2-05-1 p
5. --- యువకుల్లో ఒకడు తన \VUgras{కోటును} \textsuperscript{(10)} జాగ్రత్తగా మడిచాడు. దాన్ని ఒక
తువ్వాలుపై పెడుతున్నాడు, అందులో అతని పక్కవాడు త్వరలో రెండు \VUgras{బట్టలనూ} చుడతాడు.

%%K t03-c2-06-1 p
6. --- ఆరో కుర్రాడు ఆలస్యం చేస్తున్నాడు. అతని తలపై ఇంకా \VUgras{కాస్కెట్}
\textsuperscript{(11)} ఉంది. అతను తన \VUgras{చొక్కా} \textsuperscript{(12)} కానీ, తన
\VUgras{కాలరు} \textsuperscript{(13)} కానీ, తన \VUgras{చేతికట్లు} \textsuperscript{(14)} కానీ
విప్పలేదు. చేతిలో తన \VUgras{చొక్కాపైది} \textsuperscript{(17)} పట్టుకుని, ఇంకా తన
\VUgras{ప్యాంటు} \textsuperscript{(16)} మరియు \VUgras{సస్పెండర్లు} \textsuperscript{(15)}
వేసుకునే ఉన్నాడు. తర్వాతి పందానికి అతను సిద్ధంగా ఉండడు, ఖాయం.

%%K t03-c2-07-1 p
7. --- యువకుల దగ్గర, ఇరువైపులా చాలా \VUgras{ముళ్ళ మొక్కలున్న} \textsuperscript{(18)} ఒక చిన్న
దారి నది ఒడ్డుకు తీసుకెళ్తుంది; అక్కడ జేకబ్ తండ్రి \VUgras{రెల్లు} \textsuperscript{(19)} మధ్య
సాయంత్రం కాల్చే \VUgras{బాణసంచా} \textsuperscript{(20)} సిద్ధం చేస్తున్నాడు.

%%K t03-tit-11 sub
\VUsaut{4.4mm}
\VUpk{10.2pt}{40}{పందెం.}
\VUsaut{3.4mm}

%%K t03-c2-08-1 p
8. --- నదిపై కొందరు స్త్రీలు తమ గొడుగులతో నీడ చేసుకుంటూ \VUgras{పడవలో} \textsuperscript{(21)}
విహరిస్తున్నారు. చాలా ఆసక్తికరమైన ఆ పందాన్ని వాళ్ళు అనుసరిస్తున్నారు. ప్రతి
\VUgras{పోటీ పడవలో} \textsuperscript{(22)} నలుగురు \VUgras{తెడ్డు వేసేవాళ్ళు}
\textsuperscript{(24)} శక్తికొద్దీ తెడ్డు వేస్తున్నారు, ఈలోగా \VUgras{చుక్కాని పట్టేవాడు}
\textsuperscript{(26)} \VUgras{చుక్కానిని} \textsuperscript{(25)} పట్టుకుని ఆ సున్నితమైన పడవను
నడిపిస్తున్నాడు. \VUgras{తెడ్లు} \textsuperscript{(23)} తాళంతో నీటిని కొడుతున్నాయి, తేలికైన
పడవలు తల తిరిగే వేగంతో సాగుతున్నాయి.

%%K t03-c2-09-1 p
9. --- వంతెన స్తంభం దగ్గర కొందరు యువకులు \VUgras{నూనె స్తంభం} \textsuperscript{(28)} బహుమతి
కోసం పోటీపడుతున్నారు. వాళ్ళలో ఒకడు ఆ వంగే, జారే స్తంభంపై జాగ్రత్తగా నడుస్తూ చివర కట్టిన చిన్న
\VUgras{జెండాను} \textsuperscript{(29)} అందుకోవాలని చూస్తున్నాడు. అతని దురదృష్టవంతుడైన
\VUgras{పోటీదారు} \textsuperscript{(30)} నీళ్ళలో పడిపోయాడు. ఒడ్డుకు చేరేందుకు అతను
ఈదుతున్నాడు.

%%K t03-c2-10-1 p
10. --- ఇంకా పెద్ద యువకులు \VUgras{పడవ పోరు} \textsuperscript{(32)} చేస్తున్నారు,
\VUgras{రేవు} \textsuperscript{(33)} దగ్గర. ఈ ఆటలన్నిటినీ పెద్ద \VUgras{జనసమూహం}
\textsuperscript{(31)} చూస్తోంది — \VUgras{వంతెన} \textsuperscript{(27)} \VUgras{కటకటాలకు}
ఆనుకుని, \VUgras{ఒడ్డున} గుమిగూడి. అయితే కొందరు చూసేవాళ్ళు \VUgras{బహుమతి స్తంభం}
\textsuperscript{(45)} ఆటను చూసేందుకో, బహుమతుల \VUgras{పంపిణీకి} వస్తున్న
\VUgras{మేయరు బృందాన్ని} చూసేందుకో వెనక్కి తిరుగుతున్నారు.

%%K t03-c2-11-1 p
11. --- మొదట ఇదిగో కొందరు \VUgras{కుర్రాళ్ళు} \textsuperscript{(35)}, \VUgras{వాద్యకారుల}
\textsuperscript{(36)} ముందు నడుస్తూ; ఆ తర్వాత \VUgras{మేయరు} \textsuperscript{(37)}, ఆయన
సహాయకులు మరియు ఆ చిన్న ఊరి \VUgras{మున్సిపల్ సభ} వస్తారు. చివరగా \VUgras{అగ్నిమాపక సిబ్బంది}
\textsuperscript{(38)} నడుస్తారు.

%%K t03-tit-12 sub
\VUsaut{5.0mm}
\VUpk{10.2pt}{40}{తిరునాళ్ళు.}
\VUsaut{3.6mm}

%%K t03-c2-12-1 p
12. --- \VUgras{తిరునాళ్ళ మైదానంలో} \textsuperscript{(39)} జనం చాలా ఎక్కువ. రకరకాల
\VUgras{గుడారాలు} చూసేందుకు జనం వస్తారు : \VUgras{చెక్క గుర్రాలు} \textsuperscript{(40)},
ఇప్పటికే ప్రదర్శన మొదలైన \VUgras{సర్కస్} \textsuperscript{(41)}; \VUgras{బహిరంగ నృత్యశాల}
\textsuperscript{(42)}, అక్కడ ఊరి యువకులూ యువతులూ \VUgras{నృత్యపు} ఆనందాన్ని అనుభవిస్తారు.

%%K t03-c2-13-1 p
13. --- చివర మనకు \VUgras{క్రీడాంగణం} \textsuperscript{(44)} కనిపిస్తుంది, అక్కడ ధైర్యవంతుడైన
స్పానిష్ \VUgras{మాతదోర్} కోపంతో ఉన్న \VUgras{ఎద్దుతో} \textsuperscript{(45)} పోరాడతాడు; ఆ
తర్వాత \VUgras{జారుడు రైలు} \textsuperscript{(49)}, \VUgras{గురిచూసే స్థలం}
\textsuperscript{(46)} — అక్కడ \VUgras{తుపాకులు} వాడటమూ \VUgras{గురి} చూసి కాల్చడమూ సాధన
చేస్తారు.

%%K t03-c2-14-1 p
14. --- దగ్గరలోనే ఇదిగో \VUgras{తిరునాళ్ళ రంగస్థలం} \textsuperscript{(47)}, అక్కడ
\VUgras{హాస్యనాటకమూ} \VUgras{విషాదనాటకమూ} ఆడతారు. అగ్గా దగ్గరే \VUgras{మహాకాయుడి}
\textsuperscript{(48)} మరియు \VUgras{మరుగుజ్జు} గుడారం ఉంది. మొదటివాడి \VUgras{ఎత్తు} రెండు
మీటర్ల పదిహేను సెంటీమీటర్లు; రెండోవాడు ఒక పిల్లవాడంత కూడా ఉండడు.

%%K t03-c2-15-1 p
15. --- చివరగా ఇదిగో పెద్ద \VUgras{జంతుప్రదర్శనశాల} \textsuperscript{(50)} — దాని
\VUgras{క్రూరమృగాలతో}, బలమైన \VUgras{గోళ్ళున్న} దాని \VUgras{పులితో} \textsuperscript{(51)},
దట్టమైన \VUgras{జూలున్న} దాని \VUgras{సింహంతో} \textsuperscript{(52)}, దాని రెండు
\VUgras{జిరాఫీలతో} \textsuperscript{(53)} మరియు తన \VUgras{తొండాన్ని} ఎంతో నేర్పుగా వాడే దాని
\VUgras{ఏనుగుతో} \textsuperscript{(54)}.

%%K t03-c2-16-1 p
16. --- ఆ జంతువులకు శిక్షణ ఇచ్చే \VUgras{మచ్చిక చేసేవాడు} \textsuperscript{(55)} గుడారం తలుపు
దగ్గర ఉన్నాడు.

\VUsaut{6.0mm}
\VUfilet{22mm}
